
\documentclass{article}
%%%%%%%%%%%%%%%%%%%%%%%%%%%%%%%%%%%%%%%%%%%%%%%%%%%%%%%%%%%%%%%%%%%%%%%%%%%%%%%%%%%%%%%%%%%%%%%%%%%%%%%%%%%%%%%%%%%%%%%%%%%%%%%%%%%%%%%%%%%%%%%%%%%%%%%%%%%%%%%%%%%%%%%%%%%%%%%%%%%%%%%%%%%%%%%%%%%%%%%%%%%%%%%%%%%%%%%%%%%%%%%%%%%%%%%%%%%%%%%%%%%%%%%%%%%%
%TCIDATA{OutputFilter=LATEX.DLL}
%TCIDATA{Version=5.50.0.2953}
%TCIDATA{<META NAME="SaveForMode" CONTENT="1">}
%TCIDATA{BibliographyScheme=Manual}
%TCIDATA{Created=Monday, December 21, 2009 10:32:54}
%TCIDATA{LastRevised=Monday, December 21, 2009 17:21:45}
%TCIDATA{<META NAME="GraphicsSave" CONTENT="32">}
%TCIDATA{<META NAME="DocumentShell" CONTENT="Standard LaTeX\Blank - Standard LaTeX Article">}
%TCIDATA{CSTFile=40 LaTeX article.cst}

\newtheorem{theorem}{Theorem}
\newtheorem{acknowledgement}[theorem]{Acknowledgement}
\newtheorem{algorithm}[theorem]{Algorithm}
\newtheorem{axiom}[theorem]{Axiom}
\newtheorem{case}[theorem]{Case}
\newtheorem{claim}[theorem]{Claim}
\newtheorem{conclusion}[theorem]{Conclusion}
\newtheorem{condition}[theorem]{Condition}
\newtheorem{conjecture}[theorem]{Conjecture}
\newtheorem{corollary}[theorem]{Corollary}
\newtheorem{criterion}[theorem]{Criterion}
\newtheorem{definition}[theorem]{Definition}
\newtheorem{example}[theorem]{Example}
\newtheorem{exercise}[theorem]{Exercise}
\newtheorem{lemma}[theorem]{Lemma}
\newtheorem{notation}[theorem]{Notation}
\newtheorem{problem}[theorem]{Problem}
\newtheorem{proposition}[theorem]{Proposition}
\newtheorem{remark}[theorem]{Remark}
\newtheorem{solution}[theorem]{Solution}
\newtheorem{summary}[theorem]{Summary}
\newenvironment{proof}[1][Proof]{\noindent\textbf{#1.} }{\ \rule{0.5em}{0.5em}}
\input{tcilatex}

\begin{document}


\section{Notes on David Baron's Corporate Social Responsibility and
Entrepreneurship JEMS07}

\subsection{Introduction}

This paper, which builds on work by Graff Zivin and A. Small (2005), offers
a simple but rich general equilibrium framework within which to understand
and challenge Milton Friedman's (1970) arguments against CSR. The simplest
version of Friedman's critique is that activities by the managers of firms
amounted a costly distraction and and an  abandonment of their fiduciary
responsibility to "conduct business in accord with [shareholders'] desires,
which generally will be to make as much money as possible." \ \ One version
of this critique is that when managers engage in CSR\ activities they reduce
the NPV of cashflows that can be paid out as dividends or re-invested in the
firm, which depresses stock prices and, in effect expropriates value from
shareholders who care less about CSR activities. \ \ The subtler version of
Friedman's critique, formalized by Andreoni () and others was that even if
shareholders desired to 'do good' via CSR activities, it would tend to crowd
out other forms of charitable giving. \ \ If CSR\ were a perfect substitute
for charitable giving then every additional dollar of CSR ought to crowd out
an equivalent dollar of direct giving. \ If CSR were a costly and imperfect
substitute then CSR might both crowd out and create deadweight loss.

Friedman's seeming advice to households was to go out and first maximize the
value of their portfolios by investing in profit-maximizing firms, and then
use parts of the resulting income to invest in charity. It would seem to
follow that if households and firms followed this edict then both the size
of the economy and the size of charitable giving would be maximized and CSR\
firms would be driven out of the market.

Baron builds a framework within which to evaluate Friedman's claims and
partly challenge them. Households in the economy care about both their
private consumption and the private enjoyment they get by from giving to
"social" causes.\footnote{%
Baron greatly simplifies the analysis by treating this 'social good' as
essentially another privately enjoyed consumption good (e.g. the "warm glow"
one feels from giving to the poor, or to one's alma mater) rather than as a
'public good,' although in principle the framework could and ought to be
extended in that direction. The model does not keep track of or in anyway
analyze how the proceeds of charitable giving contribute to other people's
welfare or society's public goods. \ It is as if people are contributing to
a charity, enjoy the warm glow of giving, but those charitable funds go no
further. This is an extreme simplification but one that is useful for
identifying key tradeoffs in the household's decision.} \ The framework
gives expression to shareholders' preferences over social giving. \ All
households have the same basic preferences over consumption and 'giving' but
differ in the extent to which CSR substitutes for direct giving. \ Some
households (with low $\theta )$ find \ that a A key result is that even if
CSR is only an imperfect substitute f

In peri 

\end{document}
